\documentclass[11pt,letterpaper]{article}

\usepackage[margin=0.7in]{geometry}
\usepackage[utf8]{inputenc}
\usepackage[T1]{fontenc}
\usepackage{lmodern}
\usepackage{enumitem}
\usepackage{titlesec}
\usepackage{xcolor}
\usepackage[colorlinks=true,linkcolor=linkblue,urlcolor=linkblue,citecolor=linkblue]{hyperref}
\usepackage{parskip}

% -- Color theme --
\definecolor{accent}{HTML}{2E5090}
\definecolor{linkblue}{HTML}{1A6FB5}

% -- Section formatting --
\titleformat{\section}
  {\large\bfseries\color{accent}}
  {}{0em}{}
  [\vspace{-0.5em}\textcolor{accent}{\hrule height 0.5pt}\vspace{-0.5em}]

\titleformat{\subsection}[runin]
  {\bfseries}{}{0em}{}

% -- No page numbers --
\pagestyle{empty}

% -- Tighter lists --
\setlist[itemize]{nosep, leftmargin=1.5em, label=\textcolor{accent}{\textbullet}}

\begin{document}

% -- Header --
\begin{center}
  {\LARGE\bfseries Austin Macdonald}\\[0.3em]
  \href{mailto:asmacdo@gmail.com}{asmacdo@gmail.com} \enspace$\bullet$\enspace (336) 687-0315 \enspace$\bullet$\enspace Rock Island, IL 61201\\[0.2em]
  \href{https://github.com/asmacdo}{github.com/asmacdo} \enspace$\bullet$\enspace Source: \href{https://github.com/asmacdo/resume}{github.com/asmacdo/resume}
\end{center}

% -- Summary --
\section{Summary}

Research software engineer focused on software architecture, reproducibility, provenance, and automation for complex data and computational workflows. Translates ambiguous requirements into usable, maintainable systems across open science platforms, distributed systems, and developer tooling. Consensus builder across researchers, engineers, and stakeholders. Fully remote since 2016. Deeply integrated AI workflows---from design partner for architecture and planning to autonomous agents for well-scoped work.

% -- Professional Experience --
\section{Professional Experience}

\subsection{Research Software Engineer,} \textbf{Dartmouth College -- Center for Open Neuroscience}
\hfill\textnormal{\small2022--Present}
\begin{itemize}
  \item Designed and developed systems for reproducible scientific workflows, including execution tracking and data provenance (\href{https://github.com/con/duct}{con/duct})
  \item First author on the \href{https://github.com/myyoda/principles-paper}{STAMPED properties} paper, defining operational properties for reproducible research objects, with practical application across tools including \href{https://github.com/PennLINC/babs}{BABS} and \href{https://github.com/nipoppy/nipoppy}{nipoppy}
  \item Technical lead on the \href{https://repronim.org}{ReproNim} website redesign: aligned site architecture with pedagogical goals and maintainability, advocating for an open-source GitHub-based workflow over proprietary alternatives to enable community contributions
  \item Created \href{https://github.com/con/yolo}{con/yolo}: isolated containers to unleash Claude Code for safe agentic development
  \item Operated and maintained \href{https://github.com/dandi/dandi-hub}{DANDI Hub}, an autoscaling cloud-based JupyterHub for neuroscience data exploration and processing on AWS/Kubernetes
\end{itemize}

\subsection{Software Engineer -- Operator-SDK \& OpenShift,} \textbf{Red Hat}
\hfill\textnormal{\small2019--2022}
\begin{itemize}
  \item Led documentation overhaul across teams, establishing \href{https://operatorframework.io}{operatorframework.io}, significantly improving community usability and engagement
  \item Facilitated governance transition to a CNCF community-driven model, established regular elections, and served on the Operator Framework Steering Committee
  \item Rebuilt \href{https://github.com/operator-framework/operator-sdk}{Ansible Operator} integration as a Kubebuilder plugin, enhancing modularity and developer experience
  \item Organized and delivered multiple KubeCon presentations on Kubernetes Operators, effectively communicating complex technical concepts
\end{itemize}

\subsection{Software Engineer -- Pulp Project,} \textbf{Red Hat}
\hfill\textnormal{\small2014--2019}
\begin{itemize}
  \item Key contributor to full redesign and implementation of \href{https://github.com/pulp/pulpcore}{Pulp 3}'s REST API and plugin architecture
  \item Developed and maintained Ansible automation for deployment and developer environments, reducing barriers to usage and contribution
  \item Led Docker/podman registry Pulp plugin development, introducing container image repository versioning and rollback capabilities
  \item Mentored junior engineers and managed technical priorities balancing feature development, maintenance, and downstream support
\end{itemize}

\subsection{Developer Intern,} \textbf{Center for Open Science}
\hfill\textnormal{\small2014}
\begin{itemize}
  \item Contributed tools to the \href{https://github.com/CenterForOpenScience/osf.io}{Open Science Framework}, facilitating open practices in scientific research
\end{itemize}

% -- Publications --
\section{Selected Publications}

\begin{itemize}
  \item Zhao, C., Jarecka, D., Covitz, S., \ldots \textbf{Macdonald, A.}, \ldots Satterthwaite, T.D. (2024). \href{https://doi.org/10.1162/imag_a_00074}{A reproducible and generalizable software workflow for analysis of large-scale neuroimaging data collections using BIDS Apps}. \textit{Imaging Neuroscience}.
  \item Ioanas, H.-I., \textbf{Macdonald, A.}, Halchenko, Y.O. (2024). \href{https://doi.org/10.3389/fninf.2024.1376022}{Neuroimaging article reexecution and reproduction assessment system}. \textit{Frontiers in Neuroinformatics}.
\end{itemize}

% -- Presentations --
\section{Presentations \& Community Engagement}

\begin{itemize}
  \item \textbf{KubeCon 2022:} \href{https://www.youtube.com/watch?v=_SCzRtU5RRA}{Essential Patterns For Designing And Implementing Your Operator}
  \item \textbf{SciPy 2018:} \href{https://www.youtube.com/watch?v=5czGgUG0oXA}{Reproducible Environments for Reproducible Results}
  \item \textbf{FOSDEM/CfgMgmtCamp 2018:} Pulp 3: Ready for a Test Drive
  \item \textbf{PyCon 2018 (Poster):} Own Your Dependency Pipeline with Pulp
  \item \textbf{Podcast Interview (podcast.init, 2018)} discussing software lifecycle management and Pulp
\end{itemize}

% -- Education --
\section{Education}

\begin{itemize}
  \item \textbf{University of North Carolina at Greensboro, 2013--2014:} Computer Science prerequisites
  \item \textbf{University of Illinois at Urbana-Champaign, 2010--2012:} Master of Music Performance
  \item \textbf{University of North Carolina at Greensboro, 2005--2009:} Bachelor of Music Education
\end{itemize}

% -- Teaching & Outreach --
\section{Teaching \& Outreach}

\textbf{ChickTech Workshops (2017--2019): 3D Printers for Escher Tessellations}\\
Designed and taught a curriculum to close the gender gap in STEM by teaching symmetry, 3D modeling basics with Blender, and 3D printing, and the art of M.C. Escher

\textbf{Pisgah Astronomical Research Institute (2015--2016): 3D Planets}\\
Taught STEM outreach programs for girls across North Carolina, who each created accurate tactile planetary models with Braille for local museum exhibits, promoting diverse representation in science.

\textbf{Meredith College (2012--2014): Adjunct Professor of Music}\\
Undergraduate horn lessons and mentorship.

% -- Plot Twists --
\section{Plot Twists}

\begin{itemize}
  \item Professional French Horn Player (2008--Present)
  \item Assistant French Horn Builder at Medlin Horns
\end{itemize}

% -- Skills --
\section{Skills \& Tools}

\begin{itemize}
  \item \textbf{Core:} Python, software architecture, workflow automation, open source community building, reproducibility
  \item \textbf{Data \& Workflows:} DataLad, execution tracking, provenance
  \item \textbf{Infrastructure:} Docker, Podman, Kubernetes, AWS
  \item \textbf{Tools:} GitHub Actions, Ansible, REST APIs
\end{itemize}

\end{document}
